\section*{Chapter \ref{ch:g12:secure-communication} --- Cryptography and Secure Communication}

\begin{solution}{exo:g12:secure-communication:1}
$10110010 \oplus 01100111 = 11010101$: a $1$ wherever the two bits differ. In hexadecimal, \code{b2} xor \code{67}
gives \code{d5}, that is $213$.
\end{solution}

\begin{solution}{exo:g12:secure-communication:2}
$64$ in both cases: a \omterm{def:g12:secure-communication:hash}{digest} has $256$ bits, four per hexadecimal digit, whatever the length of the message. The
\omterm{def:g12:secure-communication:hash}{digest} of the empty message starts \code{e3b0c44298fc1c14}.
\end{solution}

\begin{solution}{exo:g12:secure-communication:3}
$6 = 110$ in binary. $5^2 = 25$, which is $2$ modulo $23$; $5^4 = 2^2 = 4$; $5^6 = 5^4 \times 5^2 = 4 \times 2 = 8$.
So $5^6 \bmod 23 = 8$, in three multiplications.
\end{solution}

\begin{solution}{exo:g12:secure-communication:4}
$855^{2753} \bmod 3233 = 123$: the message $123$, whose encryption $123^{17} \bmod 3233$ is indeed $855$. The tests
compute both.
\end{solution}

\begin{solution}{exo:g12:secure-communication:5}
\smallskip\noindent\hspace*{2em}\code{def xor_bytes(a, b):}

\noindent\hspace*{4em}\code{if len(a) != len(b):}

\noindent\hspace*{6em}\code{raise ValueError("the two byte strings must have the same length")}

\noindent\hspace*{4em}\code{return bytes(x ^ y for x, y in zip(a, b))}

\smallskip\noindent Without the check, \code{zip} would stop at the shorter one, and the end of a message would be
silently lost.
\end{solution}

\begin{solution}{exo:g12:secure-communication:6}
Every message is xored with the same bytes, so the xor of any two encrypted messages is the xor of the two plain
ones, as in \code{xor_reuse.py}: guessing a few words of one reveals the same places of all the others. The fix gives
each message a number used once, a \emph{nonce}, sent in clear with the message, and builds the stream from the key
and the nonce together:

\smallskip\noindent\hspace*{2em}\code{key_stream(f"{key}:{nonce}", len(message))}

\smallskip\noindent The nonce may be public; what matters is that a key and a nonce are never used together twice.
\end{solution}

\begin{solution}{exo:g12:secure-communication:7}
\code{def dh_shared(p, a, B): return pow(B, a, p)}. The number $g$ is already inside $B = g^b \bmod p$; Alice needs
only Bob's message, her own secret and the modulus.
\end{solution}

\begin{solution}{exo:g12:secure-communication:8}
The listener tries exponents until it meets the secret, which may be anywhere between $2$ and $p - 2$: on average
about $p/2$ tries, and the output shows means from $0.30\,p$ to $0.51\,p$. One more bit doubles $p$, so it doubles
the mean. Alice's fast power makes at most two multiplications per bit of her exponent, so one more bit adds at most
two.
\end{solution}

\begin{solution}{exo:g12:secure-communication:9}
The signature was made on the \omterm{def:g12:secure-communication:hash}{digest}, so the \omterm{def:g12:secure-communication:key-pair}{public key} turns it back into the \omterm{def:g12:secure-communication:hash}{digest} modulo $n$, not into the
message number: the comparison fails, and every genuine signature is refused. (By chance it would succeed about once
in $3233$ messages.) Signer and checker must apply the same function to the message. The tests check that the wrong
function refuses the signature of \code{"Pay Bea 20"}, which the right one accepts.
\end{solution}

\begin{solution}{exo:g12:secure-communication:10}
Secret: the browser's exponent $a$; the server's exponent $b$ and the site's \omterm{def:g12:secure-communication:key-pair}{private key}; the \omterm{def:g12:secure-communication:key-pair}{private keys} of the
authorities, held by them; the shared number and the session key, on both sides. Public: $p$ and $g$, the shares
$A$ and $B$, the \omterm{def:g12:secure-communication:certificate}{certificates} and every \omterm{def:g12:secure-communication:key-pair}{public key} in them, the signature, and the encrypted bytes.
\end{solution}

\begin{solution}{exo:g12:secure-communication:11}
\smallskip\noindent\hspace*{2em}\code{def encrypt_text(public, text):}

\noindent\hspace*{4em}\code{return [apply(public, byte) for byte in text.encode()]}

\smallskip\noindent For \code{"noon"} the two \code{o} give the same number, and so do the two \code{n}. Each byte
always gives the same number, so the result is a substitution of $256$ symbols, and letter frequencies break it as
in Book 1. Worse, anyone can encrypt all $256$ bytes with the \omterm{def:g12:secure-communication:key-pair}{public key} and read the table directly. \omterm{def:g12:secure-communication:key-pair}{Public-key
encryption} is applied to one large random number, such as a symmetric key, never to letters one by one.
\end{solution}

\begin{solution}{exo:g12:secure-communication:12}
Computing modulo $p$ is compatible with products: $(xy) \bmod p = ((x \bmod p)(y \bmod p)) \bmod p$. Applied $a$
times, $(x \bmod p)^a \bmod p = x^a \bmod p$. With $x = g^b$: $(g^b \bmod p)^a \bmod p = g^{ba} \bmod p$. In the same
way Bob's number is $g^{ab} \bmod p$, and $ab = ba$.
\end{solution}

\begin{solution}{exo:g12:secure-communication:13}
$ed = 17 \times 2753 = 46\,801 = 1 + 15 \times 3120$, and $3120 = 60 \times 52$. Take $x$ not a multiple of $61$.
Fermat gives $x^{60} \bmod 61 = 1$, so $x^{ed} = x \cdot (x^{60})^{15 \times 52}$ leaves the remainder of $x$ modulo
$61$. In the same way, with $x^{52} \bmod 53 = 1$, $x^{ed}$ and $x$ have the same remainder modulo $53$. So
$x^{ed} - x$ is a multiple of $61$ and of $53$, two different primes, hence of $3233$: $x^{ed} \bmod 3233 = x$ for
$x < 3233$. (It holds for the multiples too, but they need another line; the tests check all $3233$ values.)
\end{solution}

\begin{solution}{exo:g12:secure-communication:14}
\smallskip\noindent\hspace*{2em}\code{def store_password(password, rounds=600_000):}

\noindent\hspace*{4em}\code{salt = os.urandom(16)}

\noindent\hspace*{4em}\code{digest = hashlib.pbkdf2_hmac("sha256", password.encode(), salt, rounds)}

\noindent\hspace*{4em}\code{return salt, rounds, digest}

\smallskip\noindent\hspace*{2em}\code{def check_password(stored, attempt):}

\noindent\hspace*{4em}\code{salt, rounds, digest = stored}

\noindent\hspace*{4em}\code{trial = hashlib.pbkdf2_hmac("sha256", attempt.encode(), salt, rounds)}

\noindent\hspace*{4em}\code{return hmac.compare_digest(trial, digest)}

\smallskip\noindent The number of rounds is kept so that it can be raised later, as machines get faster: old entries
are still checked with their own number, and rewritten with the new one at the user's next login.
\code{hmac.compare_digest} compares the two \omterm{def:g12:secure-communication:hash}{digests} without stopping at the first difference, so that the time the
check takes says nothing about how much of a guess was right.
\end{solution}

\begin{solution}{exo:g12:secure-communication:15}
Mallory keeps Alice's $A$ and sends Bob her own $M = g^m \bmod p$; she does the same with Bob's $B$. Alice computes
$M^a$, which Mallory knows as $A^m$; Bob computes $M^b = B^m$. Mallory now shares one key with each, decrypts each
message, reads or changes it, and encrypts it again for the other: neither notices anything. She needs only to sit on
the path and to know $p$ and $g$, which are public. The \omterm{def:g12:secure-communication:certificate}{certificate} and the signature over the exchange stop her: the
server's share comes signed by a key that a trusted chain ties to the site's name, and Mallory cannot sign her own
share with that key.
\end{solution}

\begin{solution}{pb:g12:secure-communication:1}
\textbf{1.} $5, 2, 10, 4, 20, 8, 17, 16, 11, 9, 22, 18, 21, 13, 19, 3, 15, 6, 7, 12, 14, 1$: every number from $1$ to
$22$, once each.

\textbf{2.} Then $A$ can be any of the $22$ values, and the listener has $22$ candidates for $a$. With fewer values,
there are fewer secrets to try.

\textbf{3.} No: $2^{11} \bmod 23 = 1$, so the powers of $2$ repeat after $11$ values, and only $11$ different shares
are possible.

\textbf{4.} $5 = 101$ in binary: $5^2$, $5^4$, then $5^4 \times 5$: three multiplications.

\textbf{5.} $5^5$ is the fifth value of the list, $20$; $5^{12}$ the twelfth, $18$.

\textbf{6.} $18^5 \bmod 23 = 3$ and $20^{12} \bmod 23 = 3$, as the program printed.

\textbf{7.} $20$ is the fifth power: $a = 5$, found at the fifth try.

\textbf{8.} The exponent has at most $61$ bits: at most about $2 \times 61 = 122$ multiplications.

\textbf{9.} The secret is drawn evenly from $2$ to $p - 2$, and the search stops when it reaches it, so the number of
steps is spread evenly too, with a mean near $p/2$.

\textbf{10.} Each line is the mean of only $12$ secrets; with so few, the mean wanders around $p/2$.

\textbf{11.} From $426$ to $8\,238\,342$ steps, about $19\,000$ times more; $p$ grew $16\,432$ times. The two factors
agree within the wandering of question 10.

\textbf{12.} About $2^{k - 1}$ steps for a $k$-bit prime: it doubles with each bit.

\textbf{13.} \code{def expected_steps(k): return 2 ** (k - 1)}. For $24$ bits, $8\,388\,608$ against $8\,238\,342$
counted.

\textbf{14.} $8.4 \times 10^6 / 10^9$: under a hundredth of a second.

\textbf{15.} $2^{60} / 10^9 \approx 1.15 \times 10^9$ seconds, about $36.5$ years.

\textbf{16.} $2^{2047} / 10^9$ seconds is about $5.1 \times 10^{599}$ years.

\textbf{17.} At most about $2 \times 2048 = 4096$ multiplications of $2048$-bit numbers, a small fraction of a
second.

\textbf{18.} The prime is chosen large enough against the best known method, not against brute force. Faster methods
are why the prime has thousands of bits, where brute force alone would be stopped by a few hundred.

\textbf{19.} It gives a key, but not the identity of the other end: a man in the middle can run one exchange with
each side (exercise \ref{exo:g12:secure-communication:15}). \omterm{def:g12:secure-communication:certificate}{Certificates}, in the last section, answer that.

\textbf{20.} \textbf{Brute force needs about $p/2$ tries, so the listener's work doubles with each bit of the prime:
counted from $10$ to $24$ bits, it grew from $426$ to $8.2$ million steps, always between $0.30\,p$ and $0.51\,p$. At
$10^9$ tries per second, a prime of $2048$ bits would take about $5 \times 10^{599}$ years, while Alice needs about
$4096$ multiplications.}
\end{solution}
